\documentclass[12pt,a4paper]{article}
\usepackage[utf8]{inputenc}
\usepackage[T1]{fontenc}
\usepackage{amsmath,amssymb,amsthm}
\usepackage{physics}
\usepackage{geometry}
\usepackage{enumitem}
\usepackage{hyperref}
\usepackage{fancyhdr}
\usepackage{titlesec}

\geometry{margin=1in}
\pagestyle{fancy}
\fancyhf{}
\rhead{Lecture 4}
\lhead{Quantum Mechanics --- The Theoretical Minimum}
\cfoot{\thepage}

\titleformat{\section}{\large\bfseries}{}{0em}{}
\titleformat{\subsection}{\normalsize\bfseries}{}{0em}{}

\newtheorem{postulate}{Postulate}
\newtheorem{definition}{Definition}
\newtheorem{theorem}{Theorem}

\title{\textbf{Lecture 4: Principles of Quantum Mechanics} \\ \large Postulates, Time Evolution, and the Schr\"{o}dinger Equation}
\author{Based on Leonard Susskind's \textit{The Theoretical Minimum}}
\date{}

\begin{document}
\maketitle
\thispagestyle{fancy}

% ============================================================
\section{1. The Postulates of Quantum Mechanics}
% ============================================================

\begin{postulate}[Observables]
Every observable (measurable quantity) is represented by a \textbf{Hermitian operator} $L$ acting on the state space.
\end{postulate}

\begin{postulate}[Eigenvalues as outcomes]
The possible outcomes of measuring $L$ are its \textbf{eigenvalues} $\{\lambda_i\}$.
\end{postulate}

\begin{postulate}[Distinguishable states]
Physically distinguishable states---those that can be unambiguously told apart by some experiment---are represented by \textbf{orthogonal} vectors.
\end{postulate}

\begin{postulate}[Born's rule]
If the system is in state $\ket{A}$ and observable $L$ is measured, the probability of obtaining eigenvalue $\lambda$ is
\[
    \boxed{P(\lambda) = \bigl|\braket{\lambda}{A}\bigr|^2,}
\]
where $\ket{\lambda}$ is the normalized eigenvector of $L$ with eigenvalue $\lambda$.
\end{postulate}

\noindent The quantity $\braket{\lambda}{A}$ is the \textbf{probability amplitude}; its absolute square gives the probability.

\subsection{1.1 Degenerate eigenvalues}
If several eigenvectors $\ket{\lambda_1},\ket{\lambda_2},\dots$ share the same eigenvalue $\lambda$, the total probability is
\[
    P(\lambda) = \sum_k \bigl|\braket{\lambda_k}{A}\bigr|^2.
\]

\subsection{1.2 Results are real}
Since observables are Hermitian, all eigenvalues are real---consistent with the postulate that experimental outcomes are real numbers.

% ============================================================
\section{2. Time Evolution --- Classical Recap}
% ============================================================

In classical mechanics:
\begin{itemize}[nosep]
    \item States are updated deterministically from one instant to the next.
    \item \textbf{Reversibility} (the ``minus-first law''): distinguishable states remain distinguishable under evolution; information is conserved.
\end{itemize}

% ============================================================
\section{3. Quantum Time Evolution}
% ============================================================

\begin{postulate}[Time-evolution operator]
The state at time $t$ is obtained from the state at time $0$ by a \textbf{linear operator} $U(t)$:
\[
    \ket{\psi(t)} = U(t)\ket{\psi(0)}.
\]
The same $U(t)$ applies to every initial state.
\end{postulate}

\subsection{3.1 Conservation of distinguishability}
If $\braket{\phi(0)}{\psi(0)}=0$, then $\braket{\phi(t)}{\psi(t)}=0$ for all $t$.  More generally, the inner product is preserved:
\[
    \braket{\phi(t)}{\psi(t)} = \braket{\phi(0)}{\psi(0)}, \qquad \forall\; t.
\]

\subsection{3.2 Unitarity}
Inner-product preservation requires
\[
    \bra{\phi(0)}\,U^\dagger(t)\,U(t)\,\ket{\psi(0)} = \braket{\phi(0)}{\psi(0)},\qquad \forall\;\ket{\psi},\ket{\phi}.
\]
\begin{theorem}
This holds for all vectors if and only if
\[
    \boxed{U^\dagger(t)\,U(t) = \mathbf{1}.}
\]
Such an operator $U$ is called \textbf{unitary}.
\end{theorem}

% ============================================================
\section{4. Infinitesimal Time Evolution and the Hamiltonian}
% ============================================================

\subsection{4.1 Setup}
For an infinitesimal time step $\epsilon$:
\[
    U(\epsilon) = \mathbf{1} - i\epsilon\, H + \mathcal{O}(\epsilon^2),
\]
where $H$ is some operator.  The Hermitian conjugate is
\[
    U^\dagger(\epsilon) = \mathbf{1} + i\epsilon\, H^\dagger + \mathcal{O}(\epsilon^2).
\]

\subsection{4.2 Hermiticity of $H$}
Imposing unitarity $U^\dagger U = \mathbf{1}$ to first order in $\epsilon$:
\[
    \mathbf{1} + i\epsilon(H^\dagger - H) = \mathbf{1} \quad\Longrightarrow\quad H^\dagger = H.
\]
Hence $H$ is \textbf{Hermitian}.  Since Hermitian operators represent observables, $H$ is itself an observable---the \textbf{Hamiltonian}, identified with the energy of the system.

% ============================================================
\section{5. The Schr\"{o}dinger Equation}
% ============================================================

Starting from
\[
    \ket{\psi(\epsilon)} = (\mathbf{1} - i\epsilon H)\ket{\psi(0)},
\]
rearrange:
\[
    \frac{\ket{\psi(\epsilon)}-\ket{\psi(0)}}{\epsilon} = -iH\ket{\psi(0)}.
\]
Taking $\epsilon\to 0$:
\[
    \boxed{i\frac{d}{dt}\ket{\psi(t)} = H\ket{\psi(t)}.}
\]
This is the \textbf{time-dependent Schr\"{o}dinger equation} (with $\hbar=1$).

\subsection{5.1 Interpretation}
\begin{itemize}[nosep]
    \item The Schr\"{o}dinger equation governs the evolution of an \emph{undisturbed} system (no measurement taking place).
    \item It is deterministic in the state vector, but the state vector encodes \emph{probabilities}, not certainties, for measurement outcomes.
    \item It is the quantum analogue of the classical equations of motion.
\end{itemize}

\subsection{5.2 Time-translation invariance}
If $H$ does not depend explicitly on time, the system has \textbf{time-translation invariance}, and energy is conserved.  A time-dependent $H$ (e.g.\ a changing external field) breaks energy conservation.

% ============================================================
\section{6. Finite-Time Evolution}
% ============================================================

For time-independent $H$, repeated application of infinitesimal steps yields
\[
    U(t) = e^{-iHt}.
\]
This exponential is defined by its power series and is manifestly unitary since $H$ is Hermitian:
\[
    U^\dagger(t) = e^{+iHt} = U^{-1}(t).
\]

% ============================================================
\section{7. Summary}
% ============================================================

\begin{enumerate}[nosep]
    \item Four core postulates: observables are Hermitian operators; outcomes are eigenvalues; distinguishable states are orthogonal; probabilities follow Born's rule.
    \item Time evolution is generated by a linear, unitary operator $U(t)$.
    \item Unitarity is the quantum version of information conservation (reversibility).
    \item Expanding $U$ for small times reveals the Hamiltonian $H$---a Hermitian operator (the energy).
    \item The Schr\"{o}dinger equation $i\frac{d}{dt}\ket{\psi}=H\ket{\psi}$ determines how states evolve between measurements.
\end{enumerate}

\end{document}
